\section{Síntesis y ampliación}

Toda la conferencia puede condensarse en una conclusión de dos niveles. Sobre los modelos: combinan feature abstractas, ejecutan en paralelo varias rutas de cómputo dentro de un solo forward pass y permiten que los objetivos futuros restrinjan de antemano el token actual; estas capacidades explican fenómenos como la transferencia a lenguas de pocos recursos, la combinación de conocimientos médicos, la representación compartida entre idiomas, la suma y la rima, y también explican las autodescripciones no fieles, las alucinaciones y el conflicto interno en los jailbreak. Sobre el método: la sustitución dispersa, el attribution graph y la original-model intervention aportan evidencia causal local, pero el reconstruction error, la feature ambiguity, el success-case bias y la unexplained attention hacen que siga siendo un instrumento científico imperfecto.

\begin{importantbox}{Idea central en una frase}
No hay que preguntar «¿el modelo realmente piensa?», que es una pregunta no operativa; hay que preguntar, para un prompt dado, qué representaciones internas se activan, qué rutas contribuyen en paralelo, qué objetivos futuros aparecen por adelantado y cómo refutar esta explicación mediante intervenciones.
\end{importantbox}

\subsection{Conclusiones de la clase y lista de prácticas}

\small
\begin{multicols}{2}
\begin{enumerate}
  \item Un comportamiento correcto no implica que el mecanismo sea único ni estable;
  \item Una sola neuron no suele ser la mejor unidad semántica;
  \item La feature label debe explicarse junto con su efecto causal;
  \item El CLT mejora la legibilidad, pero también puede perder fidelidad;
  \item El reconstruction error debe aparecer explícitamente en el grafo y en el informe;
  \item Si la attention no está explicada, no afirmes que ya explicaste la elección de estrategia;
  \item Las representaciones abstractas pueden reutilizarse entre idiomas y tareas;
  \item Dentro de un solo forward pass puede haber cómputo paralelo a gran escala;
  \item La autodescripción del modelo no garantiza reflejar fielmente su algoritmo interno;
  \item IDK, la respuesta y la refusal pueden ser rutas en competencia;
  \item El planning requiere evidencia de «feature anticipada + intervención que cambia la trayectoria»;
  \item El monitoreo de seguridad debe abarcar el proceso de generación, no solo la entrada;
  \item Toda conclusión sobre mecanismos debe informar el prompt, el modelo y el periodo de tiempo;
  \item Los casos de fracaso forman parte de los resultados del método tanto como los casos de éxito.
\end{enumerate}
\end{multicols}
\normalsize

\enlargethispage{4\baselineskip}
\subsection{Lecturas adicionales}

\small
\noindent\textbf{Casos interactivos: }\href{https://transformer-circuits.pub/2025/attribution-graphs/biology.html}{On the Biology of a Large Language Model}, que incluye los grafos interactivos completos de razonamiento de varios pasos, planificación de poemas, multilingüismo, suma, diagnóstico médico, alucinaciones, rechazo y CoT faithfulness.

\noindent\textbf{Artículo de métodos: }\href{https://transformer-circuits.pub/2025/attribution-graphs/methods.html}{Circuit Tracing: Revealing Computational Graphs in Language Models}, que analiza en detalle el replacement model, el Cross-Layer Transcoder, el graph pruning, la evaluación y las limitaciones. \textbf{Orden de lectura: }primero, usa estas notas para construir la escala de evidencia feature--graph--intervention; después, vuelve al artículo interactivo y haz clic en los nodos uno por uno. Cuando encuentres un grafo vistoso, busca primero el reconstruction error, los resultados de las intervenciones, los casos de fracaso y las limitaciones de la attention.
\normalsize

\end{document}
